\section{Introduction}

The integration of cryptocurrency markets into the broader financial system has transformed them from a fringe technological curiosity into a meaningful component of global risk transmission. Bitcoin in particular has become a continuously traded, highly visible, sentiment-sensitive asset whose price dynamics are monitored not only by crypto-native investors but also by portfolio managers, macro traders, and policy observers. As the market matures, an increasingly important research question emerges: can information extracted from cryptocurrencies reveal something useful about future conditions in conventional markets such as equities, metals, or currency-related instruments?

This thesis addresses that question by focusing on \emph{time-varying intermarket dependency}. Instead of attempting to forecast price levels directly, the thesis models the changing relationship between Bitcoin and conventional assets through rolling correlation structures. This distinction is important. Correlation is not constant over time, especially in periods characterized by liquidity shocks, inflation scares, monetary tightening, or abrupt shifts in investor risk appetite. When such shifts occur, the strength and sign of cross-market linkages may change faster than conventional long-horizon diversification assumptions would suggest.

The practical motivation of the thesis is equally important. A portfolio manager does not necessarily need a perfect one-day-ahead return forecast in order to make a useful decision. In many cases, what matters most is early recognition that the market environment is becoming unstable, that co-movements are changing, and that the probability of a stress episode in conventional assets is increasing. If cryptocurrency data can contribute to such early warning information, then even a moderate forecasting edge can have real value as a risk-management input.

The central research problem can therefore be stated in two layers. The first layer is methodological and scientific: are time-varying dependencies between Bitcoin and conventional assets forecastable in a disciplined out-of-sample framework, and how do machine-learning models compare with classical econometric benchmarks? The second layer is practical: can those dependency forecasts be translated into a market-relevant warning signal that helps an investor identify elevated downside risk in conventional markets more quickly?

To address these questions, the thesis constructs a reproducible forecasting pipeline in Python. The pipeline collects price data, transforms them into log returns, builds rolling-correlation targets using multiple window lengths, engineers lagged and volatility-based predictors, and evaluates a range of models in a walk-forward setting. The model set includes simple persistence benchmarks, regularized linear methods, tree-based ensemble methods, XGBoost, and a DCC-GARCH benchmark implemented in a leakage-safe way. On top of the dependency-forecasting layer, an investor-oriented signal layer is added to predict stress regimes for conventional assets on the basis of crypto-derived information and forecasted dependency dynamics.

The empirical findings point to a nuanced but useful conclusion. Dependency forecasting works. The problem is not random. However, the dominant source of predictability is persistence in the dependency process itself, meaning that the strongest baseline models are often autoregressive and naive persistence variants. This is not a weakness of the thesis; it is a substantive result. It implies that machine-learning value in this setting is more incremental and comparative than revolutionary. The second-layer signal model provides moderate predictive value, especially as an auxiliary risk filter, but does not support the claim that cryptocurrency data can perfectly predict the next direction of all conventional markets.

The contribution of the thesis is therefore threefold. First, it builds a clean, reproducible framework for forecasting rolling intermarket dependency between cryptocurrencies and conventional assets. Second, it provides a structured comparison between classical econometric modeling and machine-learning approaches under walk-forward evaluation. Third, it extends the analysis from dependency forecasting to an applied investor warning layer, thereby connecting the statistical problem to a practical decision-making use case.

The remainder of the thesis is organized as follows. Section 2 reviews the theoretical background of time-varying dependence, financial contagion, machine-learning forecasting, and model-comparison methodology. Section 3 presents the dataset, target construction, feature engineering, models, and evaluation framework. Section 4 reports the empirical results for dependency forecasting and the investor signal layer. Section 5 interprets the findings, clarifies what the results do and do not imply, and discusses limitations. Section 6 concludes and proposes directions for future research.
